\section{Introduction}\label{sec:intro}

The chromatic Smith problem asks how much information about fixed points at one subgroup is forced by chromatic acyclicity at another. Fix a prime \(p\), and let \(K(i)\) denote Morava \(K\)-theory at \(p\), with \(K(0)=H\Q\). For a finite \(p\)-group \(G\), a subgroup \(H\le G\), and an integer \(n\ge0\), define \(r_n(G,H)\) to be the least nonnegative integer \(r\) such that
\begin{equation}\label{eq:def-loss}
 K(n+r)_*(\Phi^H X)=0
 \quad\Longrightarrow\quad
 K(n)_*(\Phi^G X)=0
\end{equation}
for every finite \(p\)-local genuine \(G\)-spectrum \(X\). Here \(\Phi^H X\) denotes the underlying spectrum of the geometric fixed points, and finite spectra include compact retracts.

Let \(\ell(G,H)\) be the least length of a subnormal chain
\[
 H=G_0\trianglelefteq G_1\trianglelefteq\cdots
 \trianglelefteq G_s=G
\]
whose successive quotients are cyclic. We set \(\ell(G,G)=0\). Such chains exist in finite \(p\)-groups. The cyclic chromatic Smith theorem, applied along a chain, gives \(r_n(G,H)\le\ell(G,H)\). Our result is the converse.

\begin{theorem}\label{thm:main}
For every prime \(p\), finite \(p\)-group \(G\), subgroup \(H\le G\), and integer \(n\ge0\),
\[
 \boxed{r_n(G,H)=\ell(G,H).}
\]
\end{theorem}

Thus the optimal loss is independent of the height. For every integer \(0\le r<\ell(G,H)\), there is a finite genuine \(G\)-spectrum whose geometric \(H\)-fixed points are \(K(n+r)\)-acyclic and whose geometric \(G\)-fixed points have nonzero \(K(n)\)-homology. The theorem includes rational fixed-point detection at \(n=0\).

The embedding of \(H\) matters. For example, in
\(G=C_{p^2}\times C_p\), take
\(H_1=pC_{p^2}\times0\) and \(H_2=0\times C_p\).
Both subgroups are isomorphic to \(C_p\), but
\(G/H_1\cong C_p\times C_p\) and \(G/H_2\cong C_{p^2}\).
Chains above \(H_i\) correspond to chains in these quotients.
The first quotient requires two cyclic steps, while the second
requires one. Thus \(r_n(G,H_1)=2\) and \(r_n(G,H_2)=1\)
at every target height. A cyclic quotient of order \(p^2\)
still counts only once.

\subsection*{History and significance}

The classical starting point is Smith theory: for a finite complex with
an action of a finite \(p\)-group, mod-\(p\) acyclicity passes to fixed
points. Floyd strengthened this to an inequality between total
mod-\(p\) Betti numbers \cite[Theorem~4.4]{Floyd}; the finite
\(p\)-group statement follows by iteration from prime-order actions.
Chromatic Smith theory asks which implications survive when mod-\(p\)
homology is replaced by individual Morava theories.

Building on the nonequivariant thick subcategory theorem of Hopkins
and Smith \cite{HS}, Balmer and Sanders identified all primes of the
category of finite genuine equivariant spectra. They reduced its
topology to subgroup--height inclusions and computed the case
\(C_p\), crediting Strickland's earlier work
\cite[Introduction, Theorems~4.9 and~4.14, and Section~7]{BS}.
Barthel, Hausmann, Naumann, Nikolaus, Noel, and Stapleton determined
the finite abelian \(p\)-group case: the loss for \(H\le G\) is the minimum
number of generators of \(G/H\) \cite[Theorem~2.3]{BHNNNS}.
This replaced the original order-based prediction by a rank bound;
already \(C_{p^2}\) has loss one rather than two.
The sharp elementary abelian examples use the partition-complex work
of Arone, Dwyer, Lesh, and Mahowald
\cite[Theorems~2.2 and~4.1]{BHNNNS}. Balderrama and Kuhn later gave
an elementary proof of the abelian chromatic Smith theorem, including
target height zero \cite[Theorem~1.1 and Section~3.2]{BK}.

For nonabelian groups, a general lower bound is the rank of
the largest elementary abelian quotient of \(G\) in which \(H\)
maps to the identity. This can be smaller than the cyclic-chain upper
bound. Kuhn and Lloyd proved that chromatic acyclicity implications
are equivalent to corresponding inequalities between Morava homology
dimensions, allowing representation-theoretic constructions to give
stronger lower bounds \cite[Theorem~2.2]{KL}. They obtained equality
with \(\ell\) for every subgroup pair inside a central extension
\(C_2\to G\to E\) with \(E\) elementary abelian
\cite[Theorem~2.13]{KL}. In particular, a noncentral subgroup of order
two in \(D_8\) has loss two, although the elementary-quotient lower
bound is one. Their discussion of the general problem expresses
``cautious hope'' that \(r_n(G,H)=r_+(G,H)=\ell(G,H)\), and separately
asks whether the loss is independent of height
\cite[Section~2.2]{KL}. \Cref{thm:main} establishes the proposed
equality, and hence height independence, for all finite \(p\)-groups.

The remaining obstacle is sharpness for an arbitrary embedded
subgroup. The Hopkins--Kuhn--Ravenel characters are indexed by maps
from \(\Z_p^n\), whose images are abelian \cite[Theorem~C]{HKR}.
Thus they cannot directly supply characters with full nonabelian
image. General character rings with other source groups already
appear in Kuhn's framework \cite[Section~2]{KuhnCharacters}.
Here we select a profinite source from the finite groups detected by
telescopic cochains, derive constraints on all its open subgroups by
relative integration, and use those constraints to prove the required
fixed-point vanishing estimates.

Combined with the conjugacy and \(p\)-subnormal reductions of
Balmer--Sanders, the theorem determines the remaining finite-height
prime inclusions for arbitrary finite groups
\cite[Proposition~6.11]{BS}; \Cref{cor:general-groups} gives the
precise test. Their description of the infinite-height inclusions
and of Thomason subsets then determines the thick tensor ideals
\cite[Corollaries~8.19, 10.5, and~10.6]{BS}. The sharpness witnesses
in this paper are obtained by separation of primes.

\subsection*{The proof}

For a finite \(p\)-group \(P\) and a spectrum \(D\), write
\[
 U_P(D)=\Phi^P\Bor_P(D),
\]
where \(\Bor_P(D)\) is the genuine cofree spectrum associated to the trivial \(P\)-action on \(D\). Let \(T(l)\) be a telescope at height \(l\). The first main step is the rational estimate
\begin{equation}\label{eq:intro-rational}
 \bigl(U_P(L_{T(l)}S)\bigr)_{\Q}\ne0
 \quad\Longrightarrow\quad \ell(P,1)\le l.
\end{equation}
To prove it at positive \(l\), we rationalize morphisms in the tensor category of dualizable \(T(l)\)-local spectra and split Burnside idempotents. A maximizing construction produces a pro-\(p\) group \(\Gamma\) mapping onto \(P\), with every finite section (a finite continuous quotient of a closed subgroup) generated by at most \(l\) elements. Cochains admit evaluation factors indexed by continuous homomorphisms out of \(\Gamma\). Height zero is handled separately by rational averaging.

The key additional information comes from relative integration over
finite groupoids. Strickland's Frobenius formulation of HKR character
theory already expresses such integration through automorphism-group
weights \cite[Section~10 and Theorem~10.2]{StricklandDuality}.
We derive the weights for the source \(\Gamma\) just constructed.
Applying them to groupoids of labeled finite sets, coinduced from an
open subgroup, gives congruences
\[
 \sum_{\substack{V\le_o U\\{}[U:V]=p^h}}
 \bigl(|X^V|-|X|\bigr)
 \in p^{h+1}\Z_{(p)}
\]
for every open subgroup \(U\le\Gamma\), every finite continuous \(U\)-set \(X\), and every \(h\ge0\). The same construction gives
\(\lvert\Hom_{\cts}(\Gamma,F)\rvert/\lvert F\rvert\in\Z_{(p)}\)
for every finite group \(F\). Together with the finite-section bound,
these arithmetic conditions force a closed subnormal series with
\(\Z_p\)-factors and length at most \(l\).
The purely profinite argument uses shrinking semidirect products and an explicit \(p\)-adic limit of subgroup counts; it may be useful independently of the chromatic application.

Next, ring nilpotence and the generalized characters of Hopkins, Kuhn, and Ravenel \cite{HS,HKR} upgrade \eqref{eq:intro-rational} to
\begin{equation}\label{eq:intro-height}
 T(k)\wedge U_P(L_{T(l)}S)\ne0
 \quad\Longrightarrow\quad k+\ell(P,1)\le l.
\end{equation}
Chromatic ambidexterity and its torsion-nilpotence consequence \cite{CSY} allow this step to use arbitrary nonzero commutative local algebras. A product with increasingly large cyclic \(p\)-groups separates the height \(k\) from the cyclic length.

Finally, finite chromatic fracture converts \eqref{eq:intro-height}
into vanishing estimates for the right adjoints to subgroup geometric
fixed points. The support argument adapts the strategy of
Balmer--Sanders \cite[proof of Theorem~9.1, \textup{(B)}$\Rightarrow$\textup{(A)}]{BS}.
At a least height where the proposed bound fails, the estimates
isolate the offending subgroup component from the other components.
In the finite-height tensor-triangular spectrum, this forces the
support of a compact object to split into two nonempty disjoint
closed parts, although that support is the closure of a point and
hence irreducible. The contradiction supplies the lower bound for
every embedded subgroup.

Published inputs are recalled in \Cref{sec:background}; the character construction, congruences, profinite theorem, higher-height estimate, and support argument occupy \Cref{sec:characters,sec:congruences,sec:profinite,sec:gluing,sec:support}, respectively.
